% ==========================================================================
%  Chapter 4 — UML Architecture and Design
% ==========================================================================
\chapter{UML Architecture and Design}
\label{ch:architecture}

\epigraph{%
  Every architecture port is a negotiation between the abstractions the kernel
  expects and the mechanisms the hardware provides.  In UML, the ``hardware''
  is another kernel.%
}{Author's observation}

User-Mode Linux is, at its core, an architecture port.  It implements the same
interfaces that \texttt{arch/x86}, \texttt{arch/arm64}, or \texttt{arch/riscv}
implement---trap dispatch, context switching, page table management, timer
infrastructure, SMP synchronization---but the target ``hardware'' is not a CPU
in ring~0.  It is the host Linux kernel, accessed through system calls.  This
chapter describes the internal architecture of the \texttt{arch/um/} tree as it
stands after the v2 redesign: the directory layout, the three-layer abstraction
model, the backend selection mechanism, the stub process model, guest memory
management, symmetric multiprocessing, the device model, and the build system.


% --------------------------------------------------------------------------
\section{The \texttt{arch/um/} Directory}
\label{sec:arch-um-dir}

In the Linux source tree, every supported architecture lives under a dedicated
subdirectory of \texttt{arch/}.  UML is no exception: \texttt{arch/um/}
contains roughly 40{,}000 lines of C (excluding tests and the archived v1 KVM
backend) spread across six major subdirectories and their associated headers.
What makes UML special is that the ``hardware abstraction layer'' at the bottom
of this code does not talk to PCI buses, interrupt controllers, or MMU hardware.
Instead it calls \syscall{clone}, \syscall{mmap}, \syscall{mprotect},
\syscall{futex}, \syscall{ioctl}, and \syscall{sigaction} on the host
kernel~\cite{dike2000,dike2006book}.

\subsection{Directory Structure}

Table~\ref{tab:arch-um-dirs} summarizes the top-level layout.  Two features of
this structure deserve emphasis.

\begin{table}[ht]
\centering
\caption{Top-level directories under \texttt{arch/um/}.}
\label{tab:arch-um-dirs}
\rowcolors{2}{UMLPanel}{white}
\begin{tabularx}{\textwidth}{l Y}
\toprule
\textbf{Directory} & \textbf{Purpose} \\
\midrule
\texttt{kernel/}   & Core kernel-side logic: trap dispatch (\texttt{trap.c}),
                      TLB synchronization (\texttt{tlb.c}), SMP
                      (\texttt{smp.c}), timekeeping (\texttt{time.c}), backend
                      arbiter (\texttt{backend.c}), Layer~2 static-key hooks
                      (\texttt{hooks.c}, \texttt{hooks\_bench.c}), template
                      pause (\texttt{template\_pause.c}), process management
                      (\texttt{process.c}), and the linker script
                      (\texttt{vmlinux.lds.S}). \\
\texttt{os-Linux/} & Host-side userspace code: signal handling
                      (\texttt{signal.c}), memory operations (\texttt{mem.c}),
                      process management (\texttt{process.c}), SMP pthread
                      lifecycle (\texttt{smp.c}), startup and probe
                      (\texttt{start\_up.c}), and the \texttt{skas/}
                      subdirectory containing the SKAS trap loop
                      (\texttt{process.c}) and stub syscall batching
                      (\texttt{mem.c}). These files are compiled with
                      \texttt{USER\_CFLAGS}---they link against glibc and the
                      host kernel's UAPI. \\
\texttt{backend/}  & Pluggable backend implementations.  Currently:
                      \texttt{seccomp/} (the mainline seccomp-SIGSYS backend),
                      \texttt{kvm-v2/} (the v2 KVM backend), \texttt{common/}
                      (shared backend utilities), \texttt{contract/} (KUnit
                      conformance tests), and \texttt{kvm-v1-archive/} (the
                      archived v1 implementation, gated by
                      \texttt{depends on BROKEN}). \\
\texttt{drivers/}  & UML-specific device drivers: the \texttt{ubd} block
                      device, vector and vector2 network drivers, console/TTY
                      drivers, the management console (\texttt{mconsole}),
                      virtio-uml, vfio-uml, and the host audio passthrough. \\
\texttt{include/}  & Headers, organized into \texttt{asm/} (kernel-internal
                      architecture headers), \texttt{shared/} (headers
                      reachable from both kernel and USER translation units),
                      and \texttt{uapi/} (userspace-visible ABI). \\
\texttt{configs/}  & Default Kconfig files: \texttt{x86\_64\_defconfig},
                      \texttt{i386\_defconfig}, \texttt{base\_defconfig}, and a
                      \texttt{profiles/} subdirectory with tuned configurations
                      for production, fuzzing, research, and sandbox use
                      cases. \\
\bottomrule
\end{tabularx}
\end{table}

First, the split between \texttt{kernel/} and \texttt{os-Linux/} is an artifact
of UML's dual-world nature.  Code in \texttt{kernel/} is compiled as kernel
code: it has access to \texttt{<linux/types.h>}, \texttt{<linux/spinlock.h>},
the kernel's \texttt{printk}, and all the usual kernel APIs.  Code in
\texttt{os-Linux/} is compiled as \emph{host userspace} code: it can call
\texttt{malloc}, \texttt{pthread\_create}, and POSIX signal functions directly.
The two worlds share data structures through the headers in
\texttt{include/shared/}, but they cannot freely call each other's functions
without crossing the kernel/user boundary through well-defined interfaces (the
\texttt{os\_*} function family declared in \texttt{include/shared/os.h}).

Second, the \texttt{backend/} directory is new in the v2 redesign.  Prior to
the backend ops abstraction (Section~\ref{sec:three-layer}), the seccomp and
ptrace trap mechanisms were interleaved throughout \texttt{os-Linux/skas/} and
gated by \texttt{if (using\_seccomp)} conditionals.  The v2 work extracted each
backend into its own directory with a clean ops table, making it possible to
add the KVM backend without touching the seccomp code paths.


% --------------------------------------------------------------------------
\section{The Three-Layer Architecture}
\label{sec:three-layer}

The v2 redesign formalized UML's runtime instrumentation and trap-mechanism
abstraction into three cleanly separated layers.  This decomposition is the
central architectural contribution of the redesign effort.  Each layer addresses
a different concern, and the layers compose orthogonally: any combination of
Layer~1 backend, Layer~2 gate configuration, and Layer~3 sanitizer set produces
a valid kernel.

\begin{figure}[ht]
\centering
\begin{tikzpicture}[
  layer/.style={
    draw, rounded corners=6pt, minimum width=13cm, minimum height=1.8cm,
    font=\small, align=center, text width=12cm
  },
  arr/.style={-{Stealth[length=5pt]}, thick, color=UMLMuted},
  lbl/.style={font=\footnotesize\bfseries, color=UMLMuted}
]

% Layer 3 — top
\node[layer, fill=UMLGreen!8, draw=UMLGreen!50] (l3) at (0,6.5)
  {\textbf{Layer 3: Compile-Time Wraps}\\[2pt]
   KASAN \quad KMSAN \quad KCSAN \quad KFENCE \quad KCOV\\[1pt]
   \textit{Completely orthogonal to Layers 1 and 2}};

% Layer 2 — middle
\node[layer, fill=UMLTeal!8, draw=UMLTeal!50] (l2) at (0,3.8)
  {\textbf{Layer 2: Static-Key Hot-Path Gates}\\[2pt]
   6 gates: \texttt{trace\_syscalls}, \texttt{kcov\_enabled},
   \texttt{time\_travel\_active},\\
   \texttt{kfence\_sample}, \texttt{record\_replay},
   \texttt{perf\_dispatch}\\[1pt]
   \textit{Off: $\sim$1--2\,ns (NOP/predicted branch) \quad
   On: hook fires (tracing, coverage, record/replay)}};

% Layer 1 — bottom
\node[layer, fill=UMLBlue!8, draw=UMLBlue!50] (l1) at (0,1.0)
  {\textbf{Layer 1: Backend Ops Table}\\[2pt]
   \texttt{struct um\_backend\_ops} --- 18+ operations in 5 categories\\
   seccomp (SIGSYS + futex) \quad | \quad kvm-v2 (\texttt{/dev/kvm} ioctl)\\[1pt]
   \textit{Single-backend: direct call \quad Dynamic: indirect call
   ($\sim$5--10 cycles)}};

% Cost arrows
\draw[arr] (l1.north) -- (l2.south)
  node[midway, right=3pt, lbl] {composes};
\draw[arr] (l2.north) -- (l3.south)
  node[midway, right=3pt, lbl] {composes};

% Side annotation
\node[font=\footnotesize\itshape, color=UMLMuted, text width=3cm, align=center]
  at (8.2,3.8) {Runtime\\toggleable\\per-event};
\node[font=\footnotesize\itshape, color=UMLMuted, text width=3cm, align=center]
  at (8.2,1.0) {Build-time\\or boot-time\\selection};
\node[font=\footnotesize\itshape, color=UMLMuted, text width=3cm, align=center]
  at (8.2,6.5) {Compile-time\\Kconfig};

\end{tikzpicture}
\caption{The three-layer architecture of UML's v2 redesign.  Each layer
  composes orthogonally with the others.}
\label{fig:three-layers}
\end{figure}


\subsection{Layer 1: The Backend Ops Table}
\label{sec:layer1}

The foundation of the architecture is \texttt{struct um\_backend\_ops}, defined
in \texttt{arch/um/include/shared/backend.h}.  Every backend---seccomp,
kvm-v2, and any future backend---populates a single global instance of this
structure.  The structure carries a contract version
(\texttt{UM\_BACKEND\_CONTRACT\_VERSION}, currently~2), a backend kind
enumeration, four capability flags, and 18 function pointers organized into
five categories:

\begin{lstlisting}[language=KernelC, caption={Abbreviated
  \texttt{struct um\_backend\_ops} (from
  \texttt{arch/um/include/shared/backend.h}).},
  label={lst:backend-ops}]
struct um_backend_ops {
    const char       *name;
    enum um_backend_kind kind;
    u32              contract_version;

    /* Capability flags */
    bool  uses_stub_reaper;
    bool  has_syscall_stub_fd_map;
    bool  stub_syscall_uses_futex;
    bool  stub_child_runs_seccomp;

    /* Lifecycle and trap (4) */
    int  (*probe)(void);
    int  (*init)(const struct um_backend_args *args);
    void (*shutdown)(void);
    void (*vcpu_run)(struct uml_pt_regs *regs);         /* HOT */

    /* Memory (5) */
    int  (*mm_create)(struct mm_struct *mm);
    void (*mm_destroy)(struct mm_struct *mm);
    int  (*mm_region_added)(...);                        /* HOT */
    int  (*mm_region_removed)(...);                      /* HOT */
    int  (*mm_region_protected)(...);

    /* Scheduling (4 + 1 optional) */
    int  (*thread_create)(...);
    int  (*thread_start_idle)(...);
    void (*context_switch)(struct task_struct *prev,     /* HOT */
                           struct task_struct *next);
    int  (*ipi_send)(int cpu, int vector);
    void (*tlb_kick_others)(struct mm_struct *mm);  /* may be NULL */

    /* Time (3) */
    u64  (*read_clock_ns)(void);                        /* HOT */
    int  (*set_timer)(int cpu, u64 deadline_ns,
                      enum um_timer_mode mode);
    u64  (*read_persistent_clock_ns)(void);

    /* Debug / introspection (3) */
    void (*init_thread_regs)(...);
    int  (*read_guest_regs)(...);
    int  (*write_guest_regs)(...);
};
\end{lstlisting}

The \textbf{HOT ops}---\texttt{vcpu\_run}, \texttt{mm\_region\_added},
\texttt{mm\_region\_removed}, \texttt{context\_switch}, and
\texttt{read\_clock\_ns}---are the operations invoked on every guest system
call, page fault, context switch, or timer read.  These are validated non-NULL
at boot time by \texttt{validate\_hot\_ops()} in
\umlpath{arch/um/kernel/backend.c}; a NULL hot op triggers an immediate kernel
panic with a diagnostic message identifying the missing operation.

The dispatch macro provides the performance-critical abstraction:

\begin{lstlisting}[language=KernelC, caption={The dispatch macro resolves at
  compile time in single-backend builds.}, label={lst:dispatch-macro}]
#if defined(CONFIG_UM_BACKEND_SECCOMP_ONLY)
# define um_backend_dispatch(op, ...) seccomp_##op(__VA_ARGS__)
#else
# define um_backend_dispatch(op, ...) (um_backend->op(__VA_ARGS__))
#endif
\end{lstlisting}

In a single-backend build (\kconfig{UM\_BACKEND\_SECCOMP\_ONLY}), the macro
expands to a direct function call via token pasting---zero indirect dispatch
overhead.  In a dynamic build (\kconfig{UM\_BACKEND\_DYNAMIC}), the macro
becomes a function pointer dereference through the global
\texttt{um\_backend} pointer, costing a single well-predicted indirect call
($\sim$5--10 cycles on modern x86).


\subsection{Layer 2: Static-Key Hot-Path Gates}
\label{sec:layer2}

Layer~2 addresses a design tension that plagues every kernel instrumentation
system: the choice between \emph{fast} and \emph{observable}.  Traditional
approaches force this choice at compile time (e.g., enabling
\kconfig{FTRACE} imposes a permanent overhead) or at boot time (e.g., kernel
command-line tracing parameters).  UML's static-key gates make observability a
\emph{runtime decision per event}.

The mechanism uses Linux's jump-label infrastructure
(\texttt{<linux/jump\_label.h>}).  Each gate is a
\texttt{DEFINE\_STATIC\_KEY\_FALSE} instance.  When off, the call site contains
either a 5-byte NOP (with \texttt{HAVE\_ARCH\_JUMP\_LABEL}) or a load-compare
branch that the CPU's branch predictor resolves in under 2\,ns.  When on, the
jump-label machinery patches the NOP to a jump into the slow-path handler.

Six gates are currently defined, each controlling a family of hook sites:

\begin{enumerate}[leftmargin=2em, itemsep=2pt]
\item \texttt{um\_hook\_trace\_syscalls} --- system call entry/exit tracing,
  page fault tracing, context switch tracing, IRQ entry tracing, clock read
  tracing.
\item \texttt{um\_hook\_kcov\_enabled} --- coverage-guided fuzzing feedback
  at syscall boundaries.
\item \texttt{um\_hook\_time\_travel\_active} --- deterministic time injection
  for reproducible execution.
\item \texttt{um\_hook\_kfence\_sample} --- KFENCE heap sampling at clock ticks.
\item \texttt{um\_hook\_record\_replay} --- record/replay event capture at
  syscall, page fault, context switch, IRQ, and clock boundaries.
\item \texttt{um\_hook\_perf\_dispatch} --- performance counter dispatch at
  syscall and context switch points.
\end{enumerate}

The hook helpers are declared \texttt{\_\_always\_inline} so that the
static-branch check expands \emph{at the call site}, not inside a function
prologue.  A typical helper composes multiple gates:

\begin{lstlisting}[language=KernelC, caption={The \texttt{um\_on\_syscall\_entry}
  hook helper composes four gates into a single inline expansion.},
  label={lst:hook-helper}]
static __always_inline void um_on_syscall_entry(struct pt_regs *regs)
{
    if (static_branch_unlikely(&um_hook_trace_syscalls))
        __um_trace_syscall_entry(regs);
    if (static_branch_unlikely(&um_hook_kcov_enabled))
        __um_kcov_record_syscall(regs);
    if (static_branch_unlikely(&um_hook_record_replay))
        __um_record_event_syscall_entry(regs);
    if (static_branch_unlikely(&um_hook_perf_dispatch))
        __um_perf_syscall(regs);
}
\end{lstlisting}

Gate state is toggled at runtime through debugfs:
\texttt{/sys/kernel/debug/um/hooks/<name>} accepts writes of \texttt{1} (on)
or \texttt{0} (off).  Each gate maintains per-CPU hit counters, summed
on read through \texttt{um\_hook\_stats\_read()}, for observability without
cross-CPU cacheline contention.

An in-kernel microbenchmark
(\umlpath{arch/um/kernel/hooks\_bench.c}) measures per-gate
cost at runtime via debugfs: 10 batches of 100{,}000 iterations each per hook
site, median reported.  This is the source of truth for the cost model.


\subsection{Layer 3: Compile-Time Wraps}
\label{sec:layer3}

Layer~3 comprises the standard kernel sanitizer and coverage tools that UML
supports by virtue of being a ``normal'' architecture port compiled by a
standard compiler toolchain:

\begin{itemize}[leftmargin=2em, itemsep=2pt]
\item \textbf{KASAN} (Kernel Address Sanitizer)~\cite{kasan} --- detects
  out-of-bounds and use-after-free bugs in slab/page allocator objects.
\item \textbf{KMSAN} (Kernel Memory Sanitizer)~\cite{kmsan} --- detects
  uses of uninitialized memory.
\item \textbf{KCSAN} (Kernel Concurrency Sanitizer)~\cite{kcsan} --- detects
  data races via compile-time instrumentation of memory accesses.
\item \textbf{KFENCE} --- lightweight probabilistic heap-object guard-page
  detection, sampling at timer ticks (gated by the Layer~2
  \texttt{kfence\_sample} hook).
\item \textbf{KCOV} --- coverage-guided fuzzing feedback, collecting per-task
  coverage bitmaps for tools like syzkaller~\cite{syzkaller}.
\end{itemize}

These wraps are completely orthogonal to Layers~1 and~2.  A KASAN-instrumented
UML running the seccomp backend with all Layer~2 gates off is a valid
configuration; so is a clean (no-sanitizer) UML running the kvm-v2 backend with
all gates on.  The three layers multiply rather than interfere.


\subsection{Cost Model}
\label{sec:cost-model}

Table~\ref{tab:cost-model} shows the approximate per-event overhead for each
layer and their composition into representative profiles.  The baseline is a
``bare'' guest system call round-trip (guest $\to$ trap $\to$ kernel
$\to$ return).  All figures assume x86-64, seccomp backend, single-backend
inline dispatch.

\begin{table}[ht]
\centering
\caption{Per-syscall-roundtrip cost model across layers and profiles.
  Figures are representative; actual measurements vary with host
  microarchitecture and workload.}
\label{tab:cost-model}
\rowcolors{2}{UMLPanel}{white}
\begin{tabular}{l r r r r}
\toprule
\textbf{Profile} & \textbf{Layer 1} & \textbf{Layer 2} & \textbf{Layer 3}
  & \textbf{Total overhead} \\
\midrule
\texttt{prod-fast}
  & 0 (inline)     & $\sim$4--8\,ns (4 NOPs)   & 0           & $\sim$4--8\,ns \\
\texttt{prod-with-hooks}
  & 0 (inline)     & $\sim$50--100\,ns (gates on) & 0         & $\sim$50--100\,ns \\
\texttt{research}
  & $\sim$5--10\,ns & $\sim$50--100\,ns        & 0           & $\sim$55--110\,ns \\
\texttt{fuzz}
  & 0 (inline)     & $\sim$20--40\,ns (kcov on) & KASAN: 2--3$\times$ & 2--3$\times$ \\
\texttt{fuzz-deep}
  & 0 (inline)     & $\sim$50--100\,ns          & KASAN+KMSAN: 4--6$\times$ & 4--6$\times$ \\
\texttt{race}
  & $\sim$5--10\,ns & $\sim$50--100\,ns        & KCSAN: 3--5$\times$ & 3--5$\times$ \\
\bottomrule
\end{tabular}
\end{table}

The key insight is that Layers~1 and~2 add \emph{nanoseconds} of overhead per
event, while Layer~3 (sanitizers) dominates total cost when enabled.  The
three-layer decomposition ensures that the fast path pays only for what it
uses, and the decision to enable observability never requires a rebuild.


% --------------------------------------------------------------------------
\section{Backend Selection Mechanism}
\label{sec:backend-selection}

Backend selection operates at two levels: \emph{build-time} configuration
through Kconfig, and \emph{boot-time} resolution through the
\texttt{backend=} kernel command-line parameter.

\subsection{Build-Time: Kconfig}

The Kconfig choice block in \umlpath{arch/um/Kconfig} offers two dispatch
modes:

\begin{description}[leftmargin=2em, style=nextline]
\item[\kconfig{UM\_BACKEND\_SECCOMP\_ONLY}]
  Single-backend inline dispatch.  The \texttt{um\_backend\_dispatch()} macro
  expands to direct \texttt{seccomp\_<op>()} calls via token pasting.  Zero
  indirect-dispatch cost.  Smallest binary.  This is the recommended mode for
  production and sandbox deployments.

\item[\kconfig{UM\_BACKEND\_DYNAMIC}]
  Multi-backend indirect dispatch.  Compiles every available backend
  (currently seccomp; kvm-v2 joins this path once its Phase~A integration
  is complete).  The dispatch macro dereferences the global
  \texttt{um\_backend} pointer, costing one well-predicted indirect call per
  site.  This is the recommended mode for distribution kernels and CI matrices
  that need a single binary to support multiple backends.
\end{description}

Both modes \texttt{select UM\_BACKEND\_SECCOMP}; the seccomp backend is always
available.  The kvm-v2 backend is gated by \kconfig{UM\_BACKEND\_KVM\_V2},
which in turn depends on \texttt{EXPERT} and \kconfig{UM\_BACKEND\_DYNAMIC}.

\begin{figure}[ht]
\centering
\begin{tikzpicture}[
  box/.style={draw, rounded corners=4pt, minimum width=3.2cm,
              minimum height=1.2cm, font=\small, align=center},
  decision/.style={draw, diamond, aspect=2.2, inner sep=1pt,
                   font=\small, align=center},
  arr/.style={-{Stealth[length=5pt]}, thick},
  note/.style={font=\scriptsize\itshape, color=UMLMuted}
]

\node[box, fill=UMLPanel] (start) at (0,0) {Kconfig\\choice};

\node[decision, fill=UMLBlue!8] (kc) at (0,-2)
  {\footnotesize SECCOMP\\ONLY?};

\node[box, fill=UMLGreen!10] (inline) at (-3.5,-4.2)
  {Inline dispatch\\direct calls};
\node[box, fill=UMLTeal!10] (dyn) at (3.5,-4.2)
  {Dynamic dispatch\\indirect calls};

\node[decision, fill=UMLTeal!8] (boot) at (3.5,-6.5)
  {\footnotesize \texttt{backend=}\\param?};

\node[box, fill=UMLAmber!10] (auto) at (0.5,-8.5) {auto\\(seccomp)};
\node[box, fill=UMLAmber!10] (seccomp) at (3.5,-8.5) {seccomp};
\node[box, fill=UMLAmber!10] (kvm) at (6.5,-8.5) {kvm-v2};

\draw[arr] (start) -- (kc);
\draw[arr] (kc) -- node[left, note] {yes} (inline);
\draw[arr] (kc) -- node[right, note] {no} (dyn);
\draw[arr] (dyn) -- (boot);
\draw[arr] (boot) -- node[left, note] {auto} (auto);
\draw[arr] (boot) -- node[left, note] {seccomp} (seccomp);
\draw[arr] (boot) -- node[right, note] {force=kvm} (kvm);

\node[note, text width=3cm] at (-3.5,-5.5)
  {\texttt{seccomp\_vcpu\_run()}\\called directly};

\end{tikzpicture}
\caption{Backend selection flow: build-time Kconfig choice followed by
  boot-time parameter resolution.}
\label{fig:backend-selection}
\end{figure}


\subsection{Boot-Time: \texttt{init\_backend()}}

The function \texttt{init\_backend()} in \umlpath{arch/um/kernel/backend.c} is
called once from \texttt{setup\_arch()} during early boot.  It resolves the
backend selection, validates HOT ops, invokes the backend's \texttt{probe()}
and \texttt{init()} lifecycle ops, and sets the immutable global
\texttt{um\_backend} pointer.

In \kconfig{UM\_BACKEND\_DYNAMIC} builds, the resolution cascade in
\texttt{pick\_dynamic\_backend()} is:

\begin{enumerate}[leftmargin=2em]
\item If \texttt{backend=force=<kind>} was specified and the requested backend
  is not compiled in, the kernel panics immediately.
\item If \texttt{backend=<kind>} was specified (non-force), the requested
  backend is used if available; otherwise the kernel falls back with a
  \texttt{pr\_warn}.
\item If \texttt{backend=auto} (the default) or no parameter was given, the
  seccomp backend is selected.
\end{enumerate}

The ptrace backend was removed entirely in the v2 refactoring.  Requesting
\texttt{backend=force=ptrace} triggers a panic with a diagnostic message
directing the user to pin to v6.16 or earlier.


% --------------------------------------------------------------------------
\section{The Stub Mechanism}
\label{sec:stub-mechanism}

The seccomp backend's core abstraction is the \emph{stub child}: a minimal
host process that serves as the execution container for a single guest
\texttt{mm\_struct}.  Each guest address space gets its own stub child, created
by \texttt{mm\_create()} and destroyed by \texttt{mm\_destroy()}.  The stub
provides the host-side process context in which guest userspace code runs
and guest system calls are intercepted.

\subsection{Stub Lifecycle}

When a new guest \texttt{mm\_struct} is created (e.g., on \texttt{fork()} or
\texttt{exec()}), the seccomp backend's \texttt{mm\_create} op clones a new
host child process.  This child---the stub---is not a conventional process.  It
executes a minimal trampoline (the ``stub code'') that:

\begin{enumerate}[leftmargin=2em]
\item Installs a seccomp BPF filter~\cite{seccomp-bpf} that traps all system
  calls, delivering \texttt{SIGSYS} to the stub's signal handler instead of
  executing the syscall on the host.
\item Maps a shared communication page (the \texttt{stub\_data} structure) at a
  fixed address in its address space.
\item Enters an infinite loop: wait on a futex in \texttt{stub\_data},
  process a batch of memory-management syscalls, and signal completion.
\end{enumerate}

\subsection{The \texttt{stub\_data} Structure}

The shared \texttt{stub\_data} page (Listing~\ref{lst:stub-data}) is the
communication channel between the UML kernel (running in the host parent
process) and the stub child.  It occupies exactly one page
(\texttt{UM\_KERN\_PAGE\_SIZE}) and contains:

\begin{lstlisting}[language=KernelC, caption={The \texttt{stub\_data}
  structure (from \texttt{arch/um/include/shared/skas/stub-data.h}).},
  label={lst:stub-data}]
struct stub_data {
    long err;
    int syscall_data_len;
    struct stub_syscall syscall_data[...] __aligned(16);

    /* Shared with signal handler (seccomp mode) */
    short restart_wait;
    unsigned int futex;   /* FUTEX_IN_CHILD or FUTEX_IN_KERN */
    int signal;
    unsigned short si_offset;
    unsigned short mctx_offset;

    struct stub_data_arch arch_data;

    /* Signal handler stack */
    unsigned char sigstack[UM_KERN_PAGE_SIZE]
        __aligned(UM_KERN_PAGE_SIZE);
};
\end{lstlisting}

The \texttt{syscall\_data} array holds a batch of \texttt{mmap}/\texttt{munmap}
operations that the stub executes on behalf of the UML kernel's TLB
synchronization layer.  Batching amortizes the futex wake/wait round-trip cost
across multiple memory operations.

\subsection{Futex-Based Synchronization}

\begin{figure}[ht]
\centering
\begin{tikzpicture}[
  proc/.style={draw, rounded corners=4pt, minimum width=2.8cm,
               minimum height=0.8cm, font=\small, fill=#1},
  msg/.style={-{Stealth[length=4pt]}, thick, color=#1},
  note/.style={font=\scriptsize, color=UMLMuted},
  timeline/.style={thick, color=UMLMuted}
]

% UML kernel (parent) timeline
\node[font=\small\bfseries, color=UMLBlue] at (-3,5.5) {UML Kernel};
\draw[timeline] (-3,5) -- (-3,-1);

% Stub child timeline
\node[font=\small\bfseries, color=UMLTeal] at (3,5.5) {Stub Child};
\draw[timeline] (3,5) -- (3,-1);

% Step 1: Kernel fills syscall batch
\node[proc=UMLBlue!10] (k1) at (-3,4.2) {Fill batch};
\node[note, anchor=east] at (-4.5,4.2) {\texttt{syscall\_data[]}};

% Step 2: futex = FUTEX_IN_CHILD, wake
\draw[msg=UMLBlue] (-3,3.5) -- node[above, note]
  {\texttt{futex} $\leftarrow$ \texttt{IN\_CHILD}; wake} (3,3.5);

% Step 3: Stub processes batch
\node[proc=UMLTeal!10] (s1) at (3,2.5) {Execute batch};
\node[note, anchor=west] at (4.5,2.5) {mmap/munmap};

% Step 4: futex = FUTEX_IN_KERN, wake parent
\draw[msg=UMLTeal] (3,1.5) -- node[above, note]
  {\texttt{futex} $\leftarrow$ \texttt{IN\_KERN}; wake} (-3,1.5);

% Step 5: Kernel reads err
\node[proc=UMLBlue!10] (k2) at (-3,0.5) {Check \texttt{err}};

% Step 6: Resume guest
\node[proc=UMLBlue!10] (k3) at (-3,-0.5) {Resume guest};
\node[proc=UMLTeal!10] (s2) at (3,0.5) {Wait on futex};

\end{tikzpicture}
\caption{Futex-based synchronization between the UML kernel and the seccomp
  stub child.  The shared \texttt{stub\_data.futex} field alternates between
  \texttt{FUTEX\_IN\_CHILD} and \texttt{FUTEX\_IN\_KERN}.}
\label{fig:stub-sync}
\end{figure}

Communication uses a single \texttt{unsigned int futex} field in
\texttt{stub\_data}.  The protocol is simple
(Figure~\ref{fig:stub-sync}):

\begin{enumerate}[leftmargin=2em]
\item The UML kernel fills the \texttt{syscall\_data} array with the batch of
  memory operations to perform.
\item The kernel sets \texttt{futex = FUTEX\_IN\_CHILD} and issues a
  \texttt{FUTEX\_WAKE} system call~\cite{futex-lwn}.
\item The stub child wakes, iterates through \texttt{syscall\_data}, executes
  each mmap/munmap, and records any error in \texttt{err}.
\item The stub sets \texttt{futex = FUTEX\_IN\_KERN} and issues
  \texttt{FUTEX\_WAKE}.
\item The kernel wakes, checks \texttt{err}, and proceeds.
\end{enumerate}

This is substantially faster than the old ptrace-based mechanism, which
required a \texttt{PTRACE\_SETREGS} + \texttt{PTRACE\_CONT} + wait round-trip
per operation~\cite{ptrace-man}.  The seccomp path achieves roughly 3--4$\times$
lower syscall round-trip latency.


% --------------------------------------------------------------------------
\section{Guest Memory Management}
\label{sec:guest-mm}

UML's guest physical memory is backed by a host-side memory object: typically
an anonymous \texttt{mmap} or a \texttt{tmpfs} file.  Guest page table changes
are not self-effectuating---there is no hardware MMU to walk guest page tables.
Instead, every guest PTE modification must be translated into a host system
call (\texttt{mmap}, \texttt{munmap}, or \texttt{mprotect}) that adjusts the
stub child's host address space to match.

\subsection{The TLB Synchronization Layer}

The mm-arbiter in \umlpath{arch/um/kernel/tlb.c} implements the
\texttt{um\_tlb\_sync()} function, which is the single drain point for pending
guest page table changes.  The flow is:

\begin{enumerate}[leftmargin=2em]
\item Guest kernel code (e.g., \texttt{do\_mmap}, \texttt{do\_munmap},
  page fault handler) modifies PTEs and calls \texttt{um\_tlb\_mark\_sync()},
  which records the affected address range in per-mm
  \texttt{sync\_tlb\_range\_from}/\texttt{sync\_tlb\_range\_to} fields.
\item Before the next guest userspace entry (\texttt{vcpu\_run}),
  \texttt{um\_tlb\_sync()} walks the per-mm page table over the marked range,
  inspects \texttt{pte\_needsync()} on each leaf, and emits a
  \texttt{struct um\_memory\_region} per dirty PTE.
\item Each region is dispatched through the backend's
  \texttt{mm\_region\_added()} or \texttt{mm\_region\_removed()} op.  For the
  seccomp backend, these ops translate into stub child mmap/munmap batches.
  For kvm-v2, they translate into KVM memslot updates.
\item After a successful drain, the mm-arbiter increments the per-mm
  \texttt{tlb\_gen} counter and, if the backend provides a
  \texttt{tlb\_kick\_others} op, invokes it to notify other vCPUs.
\end{enumerate}

\subsection{SMP TLB Generation Tracking}
\label{sec:tlb-gen}

Under SMP, multiple vCPUs may be executing code from the same
\texttt{mm\_struct} concurrently.  When one vCPU modifies a page table and
drains via \texttt{um\_tlb\_sync()}, the other vCPUs must be notified so they
can flush their local view of the guest TLB.

The mechanism uses two counters per mm, defined in
\umlpath{arch/um/include/asm/mmu.h}:

\begin{lstlisting}[language=KernelC, caption={Per-mm TLB generation counters
  for SMP coherence.}, label={lst:tlb-gen}]
typedef struct mm_context {
    /* ... */
    atomic64_t tlb_gen;
    atomic64_t tlb_gen_seen_by[NR_CPUS];
} mm_context_t;
\end{lstlisting}

\begin{itemize}[leftmargin=2em]
\item \texttt{tlb\_gen}: a monotonically increasing counter, incremented after
  each successful \texttt{um\_tlb\_sync()} drain.
\item \texttt{tlb\_gen\_seen\_by[cpu]}: for each vCPU, the most recent
  \texttt{tlb\_gen} value at which that vCPU flushed its local TLB state.
\end{itemize}

On each \texttt{vcpu\_run} dispatch, the vCPU compares its
\texttt{tlb\_gen\_seen\_by[cpu]} against the current \texttt{tlb\_gen}.  On
mismatch, it performs a guest TLB flush (for kvm-v2, a CR4.PGE toggle inside
KVM\_RUN) and updates its slot.

The cross-vCPU kick is targeted: \texttt{tlb\_kick\_others()} only sends IPIs
to vCPUs whose \texttt{tlb\_gen\_seen\_by[cpu]} is stale relative to the
current mm's \texttt{tlb\_gen}.  This narrowing eliminates the IPI storm that a
broadcast approach would cause under high mm-churn workloads.

An important subtlety: \texttt{tlb\_gen\_seen\_by} is per-(mm, cpu), not
per-vCPU globally.  A previous design used a single per-vCPU
\texttt{last\_seen\_tlb\_gen} counter, but this suffered from cross-mm
contamination: dispatching mm$_B$ (gen=5) after mm$_A$ (gen=3000) wrote 5 into
the single counter, causing a return to mm$_A$ to see a spurious lag of 2995
and potentially misleading the kicker.

\subsection{Deferred Free: The mmu\_gather Integration}

A subtlety of UML's deferred TLB-sync model is that it violates the standard
\texttt{mmu\_gather} contract.  In a hardware architecture, \texttt{mmu\_gather}
collects pages during an unmap, issues a TLB flush, and then frees the pages.
In UML, the TLB flush does not happen until the next \texttt{vcpu\_run}
dispatch.  This creates a window where freed pages could be reused by the
kernel while stale TLB entries still point to them.

The fix (rooted in \texttt{mm\_context\_t.deferred\_free\_pages}) intercepts
the page-free path: instead of releasing pages immediately, they are queued
on a per-mm deferred list and freed at the start of the next
\texttt{vcpu\_run}, after the guest TLB has been flushed.


% --------------------------------------------------------------------------
\section{SMP Implementation}
\label{sec:smp}

UML implements symmetric multiprocessing by running one host pthread per
virtual CPU.  This is a cooperative model: the host kernel schedules the
pthreads, and UML's scheduler runs within each pthread to multiplex guest
tasks.

\subsection{Thread Model}

\begin{figure}[ht]
\centering
\begin{tikzpicture}[
  cpu/.style={draw, rounded corners=6pt, minimum width=3cm,
              minimum height=4cm, fill=#1, font=\small},
  task/.style={draw, rounded corners=3pt, minimum width=2.2cm,
               minimum height=0.6cm, fill=white, font=\scriptsize},
  arr/.style={-{Stealth[length=4pt]}, thick, color=UMLMuted},
  lbl/.style={font=\footnotesize\bfseries}
]

% Host process box
\node[draw, dashed, rounded corners=8pt, minimum width=14cm,
      minimum height=6.5cm, fill=UMLPanel, label={[lbl]above:Host Process
      (UML binary)}] at (0,0) {};

% vCPU 0
\node[cpu=UMLBlue!10] (cpu0) at (-4.5,0) {};
\node[lbl, color=UMLBlue] at (-4.5,1.6) {vCPU 0};
\node[font=\scriptsize, color=UMLMuted] at (-4.5,1.1)
  {pthread 0 (pinned)};
\node[task] at (-4.5,0.3) {guest task A};
\node[task] at (-4.5,-0.5) {guest task B};
\node[task, fill=UMLGreen!10] at (-4.5,-1.3) {idle thread};

% vCPU 1
\node[cpu=UMLBlue!10] (cpu1) at (-0.5,0) {};
\node[lbl, color=UMLBlue] at (-0.5,1.6) {vCPU 1};
\node[font=\scriptsize, color=UMLMuted] at (-0.5,1.1)
  {pthread 1 (pinned)};
\node[task] at (-0.5,0.3) {guest task C};
\node[task, fill=UMLGreen!10] at (-0.5,-0.5) {idle thread};

% vCPU N-1
\node[cpu=UMLBlue!10] (cpuN) at (4.5,0) {};
\node[lbl, color=UMLBlue] at (4.5,1.6) {vCPU $N{-}1$};
\node[font=\scriptsize, color=UMLMuted] at (4.5,1.1)
  {pthread $N{-}1$ (pinned)};
\node[task] at (4.5,0.3) {guest task D};
\node[task] at (4.5,-0.5) {guest task E};
\node[task, fill=UMLGreen!10] at (4.5,-1.3) {idle thread};

% Dots
\node[font=\Large, color=UMLMuted] at (2,0) {\ldots};

% IPI arrows
\draw[arr, color=UMLAmber, bend left=20] (-3.2,-2) to
  node[below, font=\scriptsize, color=UMLAmber] {IPI (reschedule, call, TLB)}
  (-1.8,-2);
\draw[arr, color=UMLAmber, bend left=20] (1,-2) to (3.2,-2);

\end{tikzpicture}
\caption{UML SMP thread model: one host pthread per virtual CPU, each pinned
  to a host CPU via \texttt{pthread\_setaffinity\_np}.  Guest tasks are
  scheduled within each pthread by UML's copy of the Linux scheduler.}
\label{fig:smp-model}
\end{figure}

At boot time, CPU~0 runs on the main thread.  For each additional CPU
($1$ through $N{-}1$), \texttt{os\_start\_cpu\_thread()} in
\umlpath{arch/um/os-Linux/smp.c} calls \texttt{pthread\_create()} to spawn a
new host thread.  Each thread is immediately pinned to the corresponding host
CPU via \texttt{pthread\_setaffinity\_np()} (SMP-T37):

\begin{lstlisting}[language=KernelC, caption={CPU affinity pinning for vCPU
  threads (from \texttt{arch/um/os-Linux/smp.c}).}, label={lst:affinity}]
cpu_set_t cpuset;
CPU_ZERO(&cpuset);
CPU_SET(cpu, &cpuset);
pthread_setaffinity_np(cpu_threads[cpu], sizeof(cpuset), &cpuset);
\end{lstlisting}

This pinning ensures that each guest CPU's host thread always executes on the
same physical CPU, which is critical for the kvm-v2 backend's per-host-CPU
vCPU pool: a given guest CPU always dispatches against the same KVM vCPU file
descriptor.

\subsection{Context Switching}

A guest context switch is a call to the backend's \texttt{context\_switch(prev,
next)} op.  For the seccomp backend, this involves saving and restoring the
register state of the stub child process.  For kvm-v2, it involves saving
FPU state from the per-CPU vCPU into \texttt{prev->thread} and loading
\texttt{next->thread}'s saved state.

\subsection{IPI Delivery}

Inter-processor interrupts are sent through \texttt{um\_backend\_dispatch(ipi\_send, cpu, vector)}.  Four IPI vectors are defined:

\begin{description}[leftmargin=2em, font=\ttfamily, itemsep=2pt]
\item[UML\_IPI\_RES] Reschedule request (triggers \texttt{scheduler\_ipi()}).
\item[UML\_IPI\_CALL\_SINGLE] Single-function call
  (\texttt{generic\_smp\_call\_function\_single\_interrupt()}).
\item[UML\_IPI\_CALL] Broadcast function call
  (\texttt{generic\_smp\_call\_function\_interrupt()}).
\item[UML\_IPI\_STOP] CPU offline (enter infinite pause loop).
\end{description}

The template-pause mechanism (\umlpath{arch/um/kernel/template\_pause.c})
builds on top of the IPI infrastructure: it allows the UML kernel to quiesce
all vCPUs to a named ``ready point,'' \texttt{SIGSTOP} the host process for
an external supervisor to fork, and resume with per-child identity blobs.
This is the foundation of the pool-based deployment model.

\subsection{The \texttt{migrate\_disable} Fix (SMP-T13)}
\label{sec:smp-t13}

A critical correctness bug in the kvm-v2 backend's SMP support was discovered
and fixed in the SMP-T13 work item (May~2026).  The issue: UML builds with
\kconfig{PREEMPT\_VOLUNTARY}, which makes \texttt{preempt\_disable()} a no-op.
The kvm-v2 backend's \texttt{vcpu\_run} looks up the per-host-CPU vCPU pointer
using \texttt{smp\_processor\_id()} and then enters \texttt{KVM\_RUN}.  If the
host scheduler migrates the pthread to a different host CPU between the
lookup and the ioctl, the dispatch proceeds against the wrong vCPU---writing
register state to the wrong \texttt{kvm\_run} mmap, executing the wrong
vCPU's code.

The fix uses \texttt{migrate\_disable()}, which pins the running task to its
current host CPU regardless of the preempt-count semantics.  Voluntary
scheduling (e.g., inside \texttt{handle\_syscall} via \texttt{sched\_yield})
is still allowed---the key invariant is that when the task \emph{resumes}, it
resumes on the \emph{same} CPU, so the per-CPU vCPU pointer remains valid.


% --------------------------------------------------------------------------
\section{Device Model}
\label{sec:device-model}

UML provides a set of paravirtualized devices that bridge the guest kernel's
standard driver interfaces to host resources.  These are not emulations of
physical hardware; they are purpose-built drivers that use host system calls
for I/O.

\begin{table}[ht]
\centering
\caption{UML device drivers and their host-side backing mechanisms.}
\label{tab:devices}
\rowcolors{2}{UMLPanel}{white}
\begin{tabularx}{\textwidth}{l l Y}
\toprule
\textbf{Device} & \textbf{Source} & \textbf{Description} \\
\midrule
\texttt{ubd} &
  \texttt{ubd\_kern.c} &
  Block device backed by host files (typically raw disk images or sparse
  files).  Supports multiple partitions, copy-on-write layers, and synchronous
  or asynchronous I/O modes. \\

Vector network &
  \texttt{vector\_kern.c} &
  Legacy network driver using vector (scatter-gather) transport.  Supports
  TAP, raw sockets, and various encapsulation modes. \\

Vector2 network &
  \texttt{vector2\_core.c} &
  Next-generation network driver with a modular transport/model/queue
  architecture.  Supports TAP devices, host file descriptors, ethtool
  statistics, and pluggable queue management.  Default in the v2
  redesign. \\

Console/TTY &
  \texttt{stdio\_console.c} &
  Serial-like consoles mapped to host stdio, PTYs, or named pipes.  Multiple
  virtual consoles supported (\texttt{con0}, \texttt{con1}, \ldots),
  including the stderr console for early-boot diagnostics. \\

\texttt{mconsole} &
  \texttt{mconsole\_kern.c} &
  UML-specific management console accessible via a Unix domain socket.
  Supports runtime commands: \texttt{halt}, \texttt{reboot}, \texttt{config}
  (hot-add devices), \texttt{sysrq}, \texttt{cad} (Ctrl-Alt-Del), and
  \texttt{proc} (read /proc entries from outside). \\

\texttt{virtio-uml} &
  \texttt{virtio\_uml.c} &
  Standard virtio devices using vhost-user sockets.  Allows UML to use the
  same virtio device backends (e.g., vhost-user-blk, vhost-user-net) that
  QEMU and other VMMs use. \\

\texttt{hostfs} &
  \texttt{fs/hostfs/} &
  Passthrough filesystem that maps a host directory tree directly into the
  guest's VFS namespace.  Not a block device driver but a filesystem type
  (\texttt{mount -t hostfs}).  Invaluable for development: edit files on
  the host, see them immediately inside the guest. \\

\texttt{vfio-uml} &
  \texttt{vfio\_kern.c} &
  VFIO device passthrough for assigning host PCI devices to the UML guest
  (requires IOMMU support on the host). \\

Harddog &
  \texttt{harddog\_kern.c} &
  Software watchdog timer backed by a host helper process. \\

RTC &
  \texttt{rtc\_kern.c} &
  Real-time clock providing time-of-day to the guest, backed by host
  \texttt{clock\_gettime}. \\

Random &
  \texttt{random.c} &
  Entropy source backed by host \texttt{/dev/urandom}. \\
\bottomrule
\end{tabularx}
\end{table}

The device model is deliberately thin.  Because UML runs as a userspace process,
it cannot intercept hardware interrupts or DMA transactions.  Instead, UML
devices use file descriptor-based I/O and signal-driven notification
(typically \texttt{SIGIO}) to simulate device interrupts.  The
\umlpath{arch/um/kernel/sigio.c} and \umlpath{arch/um/os-Linux/sigio.c} files
manage the translation of host fd readability events into UML interrupt
delivery.

The vector2 network driver deserves special mention: it replaces the legacy
vector driver with a cleanly layered architecture separating the transport
mechanism (how packets reach the host), the device model (how the kernel sees
the NIC), and the queue management (how packets are buffered and batched).
Each layer is independently testable via KUnit, and the transport layer is
pluggable---adding a new transport (e.g., AF\_XDP) requires implementing a
small interface without touching the core driver.


% --------------------------------------------------------------------------
\section{Build System}
\label{sec:build-system}

Building UML requires specifying \texttt{ARCH=um} on the make command line.
This selects the \umlpath{arch/um/} directory as the architecture-specific
source tree and uses the host's native compiler (no cross-compilation toolchain
needed, since the ``target'' is the host architecture itself running in
userspace).

\subsection{Basic Build}

\begin{lstlisting}[language=bash, caption={Standard UML build commands.},
  label={lst:build}]
# Configure
make ARCH=um O=build/um x86_64_defconfig

# Optional: customize
make ARCH=um O=build/um menuconfig

# Build
make ARCH=um O=build/um -j$(nproc)

# Result: build/um/vmlinux (an ELF binary)
\end{lstlisting}

The output is a standard ELF binary named \texttt{vmlinux} that can be executed
directly from the command line:

\begin{lstlisting}[language=bash]
./build/um/vmlinux mem=512M rootfstype=hostfs \
    rw init=/bin/bash
\end{lstlisting}

\subsection{The \texttt{O=} Discipline}

An important operational requirement is to \emph{always} use the \texttt{O=}
out-of-tree build directory.  A bare \texttt{make ARCH=um} (without \texttt{O=})
contaminates the source tree with generated files that opaquely break future
\texttt{O=} builds.  Recovery requires \texttt{make mrproper} (which removes
all generated files, including \texttt{.config}).  All examples in this report
and in the project's CI configurations use the \texttt{O=build/um} convention.

\subsection{Kconfig Structure}

UML's Kconfig inherits from the host architecture (\texttt{x86} on x86-64
hosts) but adds UML-specific options under the \texttt{UML}-prefixed namespace.
Key configuration groups include:

\begin{description}[leftmargin=2em, style=nextline]
\item[\kconfig{UM\_BACKEND\_*}] Backend selection (Section~\ref{sec:backend-selection}).
\item[\kconfig{UM\_BACKEND\_KVM\_V2}] The v2 KVM backend (requires \texttt{EXPERT}).
\item[\kconfig{UM\_WORKER\_PROCESS}] Per-mm host worker process model.
\item[\kconfig{UM\_TEMPLATE\_PAUSE\_FORK}] Fork-on-resume pool model.
\item[\kconfig{UML\_NET\_*}] Network transport selection (TAP, raw, etc.).
\item[\kconfig{BLK\_DEV\_UBD}] The \texttt{ubd} block device.
\item[\kconfig{HOSTFS}] Host filesystem passthrough.
\end{description}

The \texttt{profiles/} subdirectory under \umlpath{arch/um/configs/} provides
ready-made configurations for common use cases: \texttt{prod-fast} (minimal
overhead, single-backend inline dispatch, no sanitizers),
\texttt{fuzz} (KASAN + KCOV + snapshot/forkserver), \texttt{research}
(dynamic dispatch with all hooks available), and others.

\subsection{SUBARCH and Host Architecture}

UML's build system uses a \texttt{SUBARCH} variable to determine the host
architecture.  On x86-64 systems, \texttt{SUBARCH=x86} is detected
automatically.  The host architecture's headers (e.g., register layouts,
calling conventions, signal frame structures) are pulled in from
\texttt{arch/x86/} via the \texttt{HEADER\_ARCH} mechanism in the top-level
\umlpath{arch/um/Makefile}:

\begin{lstlisting}[language=make, caption={Host architecture header
  selection.}]
HEADER_ARCH := $(SUBARCH)
ifneq ($(filter $(SUBARCH),x86 x86_64 i386),)
    HEADER_ARCH := x86
endif
HOST_DIR := arch/$(HEADER_ARCH)
\end{lstlisting}

This dual-architecture nature---\texttt{arch/um/} for UML-specific code,
\texttt{arch/x86/} for host-specific data structures---is what allows UML to
run unmodified x86-64 binaries in the guest while itself being a fundamentally
different ``architecture'' from the kernel's perspective.


% --------------------------------------------------------------------------
\section{Chapter Summary}
\label{sec:arch-summary}

UML's architecture is the product of a sustained negotiation between the
abstractions the Linux kernel expects from a hardware platform and the
mechanisms available through the host kernel's system call interface.  The v2
redesign formalized this negotiation into three orthogonal layers:

\begin{enumerate}[leftmargin=2em]
\item A \textbf{backend ops table} that encapsulates the trap mechanism behind
  a stable 18-operation interface, enabling both seccomp and KVM backends to
  coexist with zero code duplication in the shared kernel paths.
\item \textbf{Static-key gates} that make observability a per-event runtime
  decision, with nanosecond-scale overhead when disabled and no rebuild
  required to toggle.
\item \textbf{Compile-time sanitizer wraps} that compose with the first two
  layers without interference, enabling the full spectrum from bare-metal-speed
  production to deep sanitizer coverage.
\end{enumerate}

The stub mechanism, TLB generation tracking, SMP thread pinning, and device
model complete the picture of a virtual machine monitor that is simultaneously
a kernel architecture port, a testing platform, and a research vehicle.
Chapter~\ref{ch:implementation} explores the implementation details of each
subsystem, and Chapter~\ref{ch:kvm} describes how the kvm-v2 backend
leverages this architecture to add hardware-accelerated execution.
